\section{Consideraciones éticas y legales}

\paragraph{Ética, moral y ley en este proyecto.}
\begin{itemize}
	\item \textbf{La ley} fija mínimos exigibles. Aquí aplican la LFPDPPP (DOF 20-03-2025), la Ley General de los Derechos de Niñas, Niños y Adolescentes (LGDNNA) y el Código Penal Federal en materia de acceso ilícito a sistemas (arts. 211 bis 1 a 211 bis 7).
	\item \textbf{La moral} son las convicciones de la comunidad de la estancia. Por ejemplo, que las fotos de un niño no circulen fuera del centro, aunque ninguna norma lo diga con esas palabras.
	\item \textbf{La ética} es la reflexión del equipo para decidir cuando la ley no alcanza o cuando dos valores chocan, como la transparencia frente a la privacidad.
\end{itemize}
Cada criterio que salió de esas tres fuentes se convirtió en un requisito con prueba.

\subsection{Criterios éticos y legales convertidos en requisitos verificables}
{\small
\begin{longtable}{p{1.5cm}p{3.6cm}p{3.6cm}p{3.6cm}p{2.2cm}}
\toprule
\textbf{ID} & \textbf{Requisito verificable} & \textbf{Criterio de aceptación} & \textbf{Prueba o evidencia} & \textbf{Ref. matriz}\\\midrule\endhead
RF-07 & (Ley: LGDNNA art.~76) Una familia solo puede obtener datos de sus propios hijos & Pedir por id el expediente, los tutores o el historial de un niño ajeno no devuelve datos & \texttt{RF07\allowbreak{}Familia\allowbreak{}Solo\allowbreak{}Sus\allowbreak{}Hijos\allowbreak{}Test}, \texttt{Seguridad\allowbreak{}Test.\allowbreak{}S11} & Matriz, fila RF-07\\
RF-16 & (Moral y ética) No se puede etiquetar en una foto a un niño sin el permiso de imagen del uso elegido & \texttt{POST /\allowbreak{}fotos} con un niño sin permiso responde 403 y no guarda la foto & Validación en el servidor (\texttt{Rutas\allowbreak{}Contenido.\allowbreak{}kt}); PM-16 & Matriz, fila RF-16\\
RF-10 & (Ética) Cualquier familia que aparece en una foto puede borrarla para todos & La foto y su imagen dejan de existir para todas las familias & \texttt{RF10\allowbreak{}Eliminar\allowbreak{}Foto\allowbreak{}Test} & Matriz, fila RF-10\\
RNF-PRIV-01 / RF-18 & (Ley: LFPDPPP arts. 24 y 25, cancelación) Eliminar un expediente borra sus datos de punta a punta & Tras eliminarlo, no quedan fotos solo de él, imágenes, asistencias, incidentes, registros ligados ni la cuenta y los chats de su familia & \texttt{ARCOEliminacion\allowbreak{}Test} (3) & Matriz, filas RF-18 y RNF-PRIV-01\\
RNF-SEG-06 & (Ética: minimización, LFPDPPP art.~5) Ninguna foto conserva metadatos de ubicación & El archivo guardado no contiene EXIF, GPS ni el marcador APP1 & \texttt{SEGArchivos\allowbreak{}Test.\allowbreak{}la\allowbreak{}Foto\allowbreak{}Se\allowbreak{}Guarda\allowbreak{}Sin\allowbreak{}Exif\allowbreak{}Ni\allowbreak{}Gps} & Matriz, fila RNF-SEG-06\\
RNF-SEG-01 & (Ley: LFPDPPP art.~18, medidas de seguridad) Sin secretos en texto plano & La base no contiene contraseñas ni tokens legibles & \texttt{Seguridad\allowbreak{}Test.\allowbreak{}S06} & Matriz, fila RNF-SEG-01\\
RF-14 & (Ética: rendición de cuentas) Toda captura, consulta de foto o exportación queda registrada y solo la lee la Dirección & Otros roles reciben 403 al leer el registro & \texttt{RF14\allowbreak{}Registro\allowbreak{}Actividad\allowbreak{}Test} & Matriz, fila RF-14\\
RNF-SEG-07 & (Ética: integridad del expediente) Nadie puede fechar un registro en otro día cambiando la hora del teléfono & Fechas de 2001 o 2099 enviadas por el cliente se reemplazan por la del servidor & \texttt{SEGReloj\allowbreak{}Servidor\allowbreak{}Test} & Matriz, fila RNF-SEG-07\\
\bottomrule
\end{longtable}}

\subsection{Aviso de privacidad}
El aviso integral vive en la app (Perfil $\rightarrow$ \emph{Acerca de y aviso de privacidad}) y en el repositorio, en \texttt{app/\allowbreak{}src/\allowbreak{}main/\allowbreak{}assets/\allowbreak{}aviso-privacidad.\allowbreak{}md}. Los dos son el mismo archivo. Se transcribe a continuación.

\begin{itemize}
	\item \textbf{Ley y artículos a los que da cumplimiento:} Ley Federal de Protección de Datos Personales en Posesión de los Particulares, publicada en el DOF el 20 de marzo de 2025 y vigente desde el 21 de marzo de 2025: art.~2 fr.~VI (datos sensibles), art.~7 (consentimiento y su revocación), art.~8 (consentimiento expreso y por escrito para datos sensibles), arts.~14--17 (aviso de privacidad, su contenido mínimo y modalidades), art.~18 (medidas de seguridad), art.~21 (representante legal), arts.~22--26 (derechos ARCO), arts.~28, 31 y 34 (solicitud, plazos de 20 y 15 días hábiles y gratuidad), art.~36 (transferencias sin consentimiento) y arts.~40 y 42 (queja ante la Secretaría Anticorrupción y Buen Gobierno).
	\item \textbf{Normas complementarias para menores:} Reglamento de la LFPDPPP (DOF 21-12-2011, art.~89); Código Civil Federal, arts.~23 y 425 (representación por quien ejerce la patria potestad); LGDNNA, arts.~2 y 76 (interés superior y derecho a la intimidad).
	\item \textbf{Fuente:} texto vigente de la Cámara de Diputados, \url{https://www.diputados.gob.mx/LeyesBiblio/pdf/LFPDPPP.pdf}.
\end{itemize}

\begin{quote}\small
\textbf{Responsable.} La Fundación Estancias Infantiles Sonrisas de Cristal, en Monterrey, Nuevo León. Cada estancia tiene una Dirección que atiende las solicitudes de privacidad de sus familias y de su personal (art.~29).

\textbf{Datos que recopilamos.}
\begin{itemize}
	\item De niñas y niños: nombre, apellidos, fecha de nacimiento y sala. Datos de salud, que son sensibles: alergias, padecimientos, medicamentos, lesiones o síntomas al llegar e incidentes. Su día en la estancia: llegada, salida, comidas, siestas, baño, ánimo y notas. Fotografías, solo con la autorización de imagen para el uso autorizado.
	\item De tutores y personas autorizadas: nombre, parentesco y teléfono; quién dejó y quién recogió al niño; mensajes enviados a la estancia.
	\item Del personal: nombre, usuario, rol, estancia y sala.
	\item De todas las cuentas: usuario, una huella irreversible de la contraseña (PBKDF2) y un registro de actividad con la hora del servidor.
\end{itemize}
No recopilamos ubicación. A las fotos se les borran los metadatos (EXIF/GPS). Tu huella digital nunca sale del teléfono.

\textbf{Para qué los usamos.}
\begin{enumerate}
	\item Registrar la asistencia y entregar al niño solo a personas autorizadas.
	\item Llevar la bitácora diaria y avisar a la familia de incidentes y de temas de salud.
	\item La comunicación entre la familia y la estancia.
	\item Compartir con cada familia las fotos donde aparece su hijo.
	\item Reportes de asistencia; a la Fundación solo le llegan totales.
	\item Proteger las cuentas y registrar quién hizo qué.
\end{enumerate}
No los usamos con fines comerciales ni para decisiones automatizadas.

\textbf{Consentimiento.} Quien ejerce la patria potestad da su consentimiento expreso y por escrito para los datos de salud al inscribir al niño (art.~8). La autorización de imagen se registra por separado en el expediente.

\textbf{Dónde se guardan.}
\begin{itemize}
	\item En la base PostgreSQL de la Fundación, en un contenedor aislado.
	\item Viajan cifrados con HTTPS (TLS).
	\item En el teléfono, la sesión se guarda cifrada con AES-256-GCM, con una llave del Android Keystore.
	\item Cada persona ve solo lo que le corresponde.
\end{itemize}

\textbf{Con quién se comparten.} Con nadie fuera de la Fundación, salvo los casos del art.~36 (por ejemplo, un requerimiento de una autoridad competente o una emergencia médica de un niño).

\textbf{Cuánto tiempo se conservan.}
\begin{itemize}
	\item El expediente, mientras el niño esté inscrito. Al darlo de baja, se borra automáticamente a los 30 días.
	\item Las fotos, hasta que se borren. De las fotos de grupo, al eliminar un expediente se quita la etiqueta del niño.
	\item Las cuentas de familia, hasta que se elimina su último expediente, junto con sus conversaciones.
	\item Las sesiones, 30 días.
	\item El registro de actividad, mientras opere la estancia. Sus entradas ligadas a un niño se borran con el expediente.
\end{itemize}

\textbf{Derechos ARCO.} Tú, o quien ejerza la patria potestad del menor (art.~21; art.~89 del Reglamento), puedes pedir a la Dirección, por escrito, el acceso, la rectificación, la cancelación o la oposición. Respondemos en 20 días hábiles y aplicamos lo procedente en los 15 siguientes (art.~31). El trámite es gratuito (art.~34). Si no estás de acuerdo con la respuesta, puedes acudir a la Secretaría Anticorrupción y Buen Gobierno (arts.~40 y 42).

\textbf{Cambios.} Se publican en la app y se anuncian en Avisos (art.~15 fr.~VI).
\end{quote}

\paragraph{Correspondencia con el comportamiento real.}
\begin{itemize}
	\item \textbf{Acceso:} la familia ve todo lo registrado de su hijo.
	\item \textbf{Rectificación:} la Dirección edita el expediente.
	\item \textbf{Cancelación:} baja seguida de eliminación, comprobada por \texttt{ARCOEliminacion\allowbreak{}Test}.
	\item \textbf{Oposición a fotos:} la Dirección cambia la autorización de imagen, y \texttt{POST /\allowbreak{}fotos} la hace cumplir.
	\item La limpieza de metadatos y el cifrado tienen sus propias pruebas (\texttt{SEGArchivos\allowbreak{}Test}, \texttt{S06}).
\end{itemize}

\subsection{Clasificación de los datos}
Niveles de sensibilidad:
\begin{itemize}
	\item \textbf{S} (sensible, art.~2 fr.~VI): su mal uso puede causar discriminación o un riesgo grave.
	\item \textbf{A} (alto): datos de menores o que permiten localizarlos o identificarlos.
	\item \textbf{M} (medio): datos de contacto o de cuenta.
	\item \textbf{B} (bajo): datos operativos.
\end{itemize}
Todo dato del servidor viaja por TLS, vive en PostgreSQL dentro de un volumen no expuesto y se entrega solo según la regla de visibilidad por rol. En la tabla, la columna \emph{Protección} indica lo que se agrega a eso.

{\small
\begin{longtable}{p{3.6cm}p{1.3cm}p{3.2cm}p{3.6cm}p{3.0cm}}
\toprule
\textbf{Dato almacenado} & \textbf{Nivel} & \textbf{Ubicación} & \textbf{Protección aplicada} & \textbf{Conservación}\\\midrule\endhead
Nombre y apellidos del niño & A & \texttt{ninos.\allowbreak{}datos} & Visible solo para su sala, su familia y su Dirección & Inscripción + 30 días de baja\\
Fecha de nacimiento del niño & A & \texttt{ninos.\allowbreak{}datos} & Igual que el nombre & Inscripción + 30 días\\
Alergias del niño & S & \texttt{ninos.\allowbreak{}datos} & Consentimiento escrito; la Fundación no lo ve & Inscripción + 30 días\\
Padecimientos y medicamentos & S & \texttt{ninos.\allowbreak{}datos} & Igual que alergias & Inscripción + 30 días\\
Permisos de imagen (interno y redes) & M & \texttt{ninos.\allowbreak{}datos} & Validados en el servidor al publicar fotos & Inscripción + 30 días\\
Estado del expediente y fecha de eliminación & B & \texttt{ninos.\allowbreak{}datos} & Solo lo cambian \texttt{/\allowbreak{}baja} y \texttt{/\allowbreak{}reactivar}, con la hora del servidor & Inscripción + 30 días\\
Personas autorizadas: nombre & A & \texttt{ninos.\allowbreak{}datos} & Visible para la sala, la familia y la Dirección & Inscripción + 30 días\\
Personas autorizadas: parentesco y teléfono & M & \texttt{ninos.\allowbreak{}datos} & Igual & Inscripción + 30 días\\
Tutor: nombre y parentesco & M & \texttt{tutores.\allowbreak{}datos} & Visible para su estancia; borrado con el último hijo & Mientras tenga un hijo inscrito\\
Tutor: teléfono & M & \texttt{tutores.\allowbreak{}datos} & Igual & Igual\\
Hora de llegada y de salida & A & \texttt{asistencias.\allowbreak{}datos} & Fecha del servidor; hora limitada a la del servidor & Inscripción + 30 días\\
Quién entregó y quién recogió & A & \texttt{asistencias.\allowbreak{}datos} & Solo personas autorizadas del expediente (RF-15) & Inscripción + 30 días\\
Lesiones, síntomas y medicamento al llegar & S & \texttt{asistencias.\allowbreak{}datos} & Visible solo para la sala, la familia y la Dirección & Inscripción + 30 días\\
Comidas, siestas, baño y ánimo & A & \texttt{eventos.\allowbreak{}datos} & Firmados por la educadora con la hora del servidor & Inscripción + 30 días\\
Notas de la educadora & A & \texttt{eventos.\allowbreak{}datos} & Igual & Inscripción + 30 días\\
Incidente: tipo y parte del cuerpo & S & \texttt{incidentes.\allowbreak{}datos} & Folio único; hora del servidor & Inscripción + 30 días\\
Incidente: descripción y acción tomada & S & \texttt{incidentes.\allowbreak{}datos} & Igual & Inscripción + 30 días\\
Imagen de la foto & A & \texttt{archivos.\allowbreak{}bytes} & Magic bytes, re-codificada sin EXIF ni GPS; descarga solo por quien ve la foto & Hasta que se borre\\
Descripción y niños etiquetados de la foto & A & \texttt{fotos.\allowbreak{}datos} & Visible solo para las familias etiquetadas & Hasta que se borre\\
Texto del mensaje & M & \texttt{mensajes.\allowbreak{}datos} & Solo participantes; no se copia al registro de actividad & Mientras exista la cuenta\\
Autor y hora del mensaje & M & \texttt{mensajes.\allowbreak{}datos} & Hora del servidor & Igual\\
Lecturas de mensajes y avisos & B & \texttt{lecturas}, \texttt{avisos\_leidos} & Por usuario & Mientras exista la cuenta\\
Quién reclamó un objeto perdido & M & \texttt{objetos.\allowbreak{}datos} & Una familia solo ve su propio reclamo & Mientras opere la estancia\\
Foto del objeto perdido & B & \texttt{archivos.\allowbreak{}bytes} & Sin metadatos; visible solo para la sala & Mientras opere la estancia\\
Usuario (login) & M & \texttt{usuarios.\allowbreak{}datos} & No visible para familias & Mientras exista la cuenta\\
Nombre y rol de la cuenta & M & \texttt{usuarios.\allowbreak{}datos} & Por estancia & Mientras exista la cuenta\\
Contraseña & S & \texttt{claves.\allowbreak{}hash} & PBKDF2-SHA256, 210\,000 iteraciones, con sal; nunca en claro & Mientras exista la cuenta\\
Token de sesión (servidor) & S & \texttt{sesiones.\allowbreak{}token\_hash} & Solo su SHA-256 & 30 días; los vencidos se purgan\\
Token de sesión (teléfono) & S & SharedPreferences de la app & AES-256-GCM con una llave del Android Keystore; \texttt{allow\allowbreak{}Backup=false} & Hasta cerrar sesión\\
Desbloqueo con huella activo & B & \texttt{huellas} & Solo un indicador; la huella no sale del teléfono & Mientras exista la cuenta\\
Bloqueo de cuenta & B & \texttt{bloqueos} & Solo la Fundación lo ve y lo quita & 3 días o hasta el desbloqueo\\
IP de intentos fallidos & M & Memoria del servidor & No se persiste & 24 horas\\
Registro de actividad & M & \texttt{registros.\allowbreak{}datos} & Solo lo lee la Dirección; hora del servidor & Mientras opere la estancia; lo ligado a un niño, con su expediente\\
\bottomrule
\end{longtable}}

\subsection{Dilemas éticos y decisiones de diseño}

\paragraph{Dilema 1: compartir fotos o proteger la imagen de los niños.}
\begin{description}
	\item[El dilema.] Las fotos tranquilizan a las familias, pero cada foto de grupo expone a otros niños cuyas familias quizá no la autorizaron. Además, el archivo puede llevar la ubicación de la estancia en sus metadatos.
	\item[Alternativas y argumentos.]
	\begin{enumerate}
		\item Compartir todas las fotos con toda la sala: máxima transparencia y menos trabajo para la educadora.
		\item No usar fotos: cero riesgo, pero se pierde un beneficio que las familias valoran.
		\item Compartir con reglas: cada foto la ven solo las familias de los niños que aparecen, con permiso de imagen por uso, y cualquiera de esas familias puede borrarla. Es más complejo, pero respeta a cada familia.
	\end{enumerate}
	\item[Postura del equipo.] La alternativa 3, por el interés superior de la niñez (LGDNNA arts.~2 y 76).
	\item[Decisión concreta.]
	\begin{itemize}
		\item Dos permisos de imagen por niño (RF-16), validados en el servidor.
		\item La visibilidad de cada foto depende de los niños etiquetados (RF-09).
		\item Cualquier familia etiquetada puede borrar la foto para todos (RF-10).
		\item Toda imagen se guarda sin EXIF ni GPS.
		\item Las fotos se toman dentro de la app y no pasan por la galería del teléfono personal.
	\end{itemize}
	\item[Dónde se observa.] Pantalla \emph{Fotos} (los niños sin permiso no se pueden etiquetar), \texttt{Rutas\allowbreak{}Contenido.\allowbreak{}kt}, \texttt{RF09\allowbreak{}Fotos\allowbreak{}Visibles\allowbreak{}Test}, \texttt{RF10\allowbreak{}Eliminar\allowbreak{}Foto\allowbreak{}Test} y \texttt{SEGArchivos\allowbreak{}Test}.
\end{description}

\paragraph{Dilema 2: trazabilidad o privacidad del registro de actividad.}
\begin{description}
	\item[El dilema.] Registrar quién hizo qué protege a los niños y permite aclarar incidentes, pero el registro mismo es un expediente sobre el personal y las familias, y puede volverse vigilancia.
	\item[Alternativas y argumentos.]
	\begin{enumerate}
		\item Registrar todo, mensajes incluidos, y que todos lo consulten: máxima transparencia.
		\item No registrar nada: máxima privacidad, pero sin rendición de cuentas.
		\item Registrar solo acciones sobre datos de niños (capturas, consultas de fotos, exportaciones, cuentas) y que solo lo lea la Dirección.
	\end{enumerate}
	\item[Postura del equipo.] La alternativa 3: rendición de cuentas proporcional (LFPDPPP art.~5, proporcionalidad).
	\item[Decisión concreta.]
	\begin{itemize}
		\item RF-14: la lectura del registro está restringida a la Dirección.
		\item El contenido de los mensajes nunca se copia al registro.
		\item La consulta de una foto por una familia sí queda registrada.
		\item Al eliminar un expediente se borran sus entradas del registro.
	\end{itemize}
	\item[Dónde se observa.] \emph{Más} $\rightarrow$ \emph{Registro de actividad} (solo Dirección), \texttt{RF14\allowbreak{}Registro\allowbreak{}Actividad\allowbreak{}Test} y \texttt{ARCOEliminacion\allowbreak{}Test}.
\end{description}

\paragraph{Dilema 3: borrar de verdad o conservar evidencia.}
\begin{description}
	\item[El dilema.] El derecho de cancelación pide borrar los datos de un niño que deja la estancia. Pero borrar al instante puede destruir evidencia de un incidente, o un error de la Dirección puede borrar un expediente activo.
	\item[Alternativas y argumentos.]
	\begin{enumerate}
		\item Borrado inmediato: cumple la cancelación al momento, pero un error es irreversible.
		\item Nunca borrar (solo desactivar): conserva la evidencia, pero contradice la ley y el aviso de privacidad.
		\item Baja con un periodo de gracia de 30 días y luego borrado completo: permite corregir errores y luego cumple la cancelación de punta a punta.
	\end{enumerate}
	\item[Postura del equipo.] La alternativa 3.
	\item[Decisión concreta.]
	\begin{itemize}
		\item El servidor rechaza eliminar un expediente activo (S20).
		\item La purga automática corre cada 6 horas.
		\item El borrado incluye las imágenes, las entradas del registro, la cuenta de la familia y sus chats (\texttt{ARCOEliminacion\allowbreak{}Test}).
		\item En la auditoría se corrigió que las imágenes quedaban huérfanas.
	\end{itemize}
	\item[Dónde se observa.] Expediente $\rightarrow$ \emph{Dar de baja} y \emph{Eliminar definitivamente}; sección 7 del aviso de privacidad.
\end{description}
